Un parent tire un traîneau sur lequel est assis son enfant. La masse totale du traîneau et de l'enfant est de \qty{\v{m}}{kg}. La corde exerce une force de \qty{\v{F}}{N} et fait un angle de \qty{\v{th}}{\degree} au-dessus de l'horizontale. Utilisez $g = \qty{9.81}{N/kg}$.

\begin{enumerate}[label=\alph*)]
    \item Déterminez la composante horizontale de la force exercée par la corde.
    \item Déterminez la composante verticale de la force exercée par la corde.
    \item Déterminez la force normale exercée par le sol sur le traîneau. Expliquez pourquoi elle est plus petite que le poids.
\end{enumerate}
